
\medskip
\phantomsection\label{par:sga3-VIA-4.5}%
\noindent\textbf{4.5.---}
We have now constructed \(F\backslash G\) and shown that it is a
direct sum of schemes of finite type over \(A\). The remaining
assertions of \hyperref[thm:sga3-VIA-3.2]{Theorem~3.2} reduce directly to assertions about the
groupoids \(\mathscr S(G^\alpha_*)\). As in Expos\'e~V, \S~6, the
second assertion in \SGARefLink{thm:sga3-VIA-3.2:item:i}{(i)} follows from the first together with \SGARefLink{thm:sga3-VIA-3.2:item:ii}{(ii)} and
\SGARefLink{thm:sga3-VIA-3.2:item:iii}{(iii)}; it is therefore enough to prove \SGARefLink{thm:sga3-VIA-3.2:item:ii}{(ii)}, (iii), and (iv)(a).

Since \(A'\) is a local \(A\)-algebra, finite and free over \(A\), the
morphism \(A\to A'\) is faithfully flat and of finite presentation.
By SGA~1, VIII (3.1, \SGARefLink{par:sga3-VIA-4.6}{4.6}, 5.4), it is therefore enough to verify the
corresponding assertions for the groupoid
\[
  \mathscr S(G^\alpha_*)\otimes_AA'.
\]
This groupoid is isomorphic to a direct sum of finitely many copies of
\(\mathscr S(G^0_*)\otimes_AA'\) (cf.\ \hyperref[par:sga3-VIA-4.4]{4.4}), so we are reduced to the
groupoid \(\mathscr S(G^0_*)\). For this last groupoid we continue to
follow the proof given in Expos\'e~V, \S~6, as we began to do in \SGARefLink{par:sga3-VIA-4.3}{4.3}.

\medskip
\phantomsection\label{par:sga3-VIA-4.6}%
\noindent\textbf{4.6.---}
Let us conclude this section with some remarks concerning \SGARefLink{lem:sga3-VIA-6.1}{Lemma~6.1}
and \S~9(a) of Expos\'e~V. Under the hypotheses and with the notation
of Expos\'e~V, \S~9(a), we seek a condition under which formation of
the cokernel of the \((\mathrm{Sch}/S)\)-groupoid \(X_*\) commutes
with a base extension
\[
  \pi:S'\longrightarrow S.
\]
The cokernels of \(X_*\) and \(X'_*\) are identified with the
cokernels of the groupoids \(U_*\) and \(U'_*\) induced by \(X_*\)
and \(X'_*\) on the quasi-sections \(U\) and \(U'\), respectively.
We are therefore reduced to the case of a
\((\mathrm{Sch}/S)\)-groupoid satisfying the hypotheses of
\SGARefLink{thm:sga3-V-4.1}{Expos\'e~V, Theorem~4.1}.

\footnote{\label{note:sga3-VIA-45}Editor's note: we have added the
following sentence.}%
Let \(Y\) be the cokernel of \(U_*\), put
\[
  Y'=Y\times_SS',
\]
and let \(Y_1\) be the cokernel of \(U'_*\); we saw in
Expos\'e~V, \S~9(a), that the canonical morphism
\[
  Y_1\longrightarrow Y'
\]
is a homeomorphism (indeed, a universal homeomorphism); we may thus
identify \(Y_1\) and \(Y'\) as topological spaces. If
\(p:U\to Y\) is the canonical morphism and
\(p':U'\to Y'\) is obtained from it by base change, we want the
following sequence of \(\mathcal O_{Y'}\)-modules to be exact:
\[
  \tag{*}\label{eq:sga3-VIA-4.6-star}
  \mathcal O_{Y'}
    \longrightarrow p'_*(\mathcal O_{U'})
    \rightrightarrows
  p'_*u'_{1*}(\mathcal O_{U'_1})
    =
  p'_*u'_{0*}(\mathcal O_{U'_1}).
\]
\footnote{\label{note:sga3-VIA-46}Editor's note: in what follows we
have modified the original; the additional hypotheses imposed there
on \(X_*\) are unnecessary.}%
Since we are under the hypotheses of \SGARefLink{thm:sga3-V-4.1}{Expos\'e~V, Theorem~4.1},
\(u_0\) and \(u_1\) are finite and locally free; moreover,
\(p\) is integral by \SGARefLink{thm:sga3-V-4.1:item:ii}{Expos\'e~V, 4.1(ii)}.
Thus \(p\) and \(p\circ u_i\) are affine, hence separated and
quasi-compact.

Consequently, if \(S'\) is flat over \(S\), EGA~III\textsubscript{1},
1.4.15 (taking account of correction Err\textsubscript{III}~25 in
EGA~III\textsubscript{2}) identifies the sequence
\eqref{eq:sga3-VIA-4.6-star} with the inverse image of the sequence
\[
  \tag{**}\label{eq:sga3-VIA-4.6-double-star}
  \mathcal O_Y
    \longrightarrow p_*(\mathcal O_U)
    \rightrightarrows
  p_*u_{1*}(\mathcal O_{U_1})
    =
  p_*u_{0*}(\mathcal O_{U_1}),
\]
which is exact.%
\footnote{\label{note:sga3-VIA-47}Editor's note: we have added both
the following sentence and the numbering \SGARefLink{prop:sga3-VIA-4.6.1}{4.6.1} below, in order to
highlight the stated result.}
An analogous argument applies when the groupoid
\(X_*\) has quasi-sections ``locally'' (cf.\ the proof of
\SGARefLink{thm:sga3-V-7.1}{Expos\'e~V, Theorem~7.1}).
We therefore obtain:

\medskip
\phantomsection\label{prop:sga3-VIA-4.6.1}%
\noindent\textbf{Proposition 4.6.1.}
\textit{Formation of the cokernel of \(X_*\) commutes with flat base
extensions when \(X_*\) has quasi-sections locally.}
